\section{Analysis}
\label{sec:analysis}

\subsection{Why TF-IDF Outperformed Mean-Pooled GloVe}

The strongest TF-IDF model exceeded the strongest GloVe model by 0.1537 macro-F1 on the 
official test set. TF-IDF preserves corpus-specific lexical importance: rare words and short 
phrases associated with sentiment receive strong weights, while common terms are discounted. 
In contrast, unweighted GloVe mean pooling compresses an entire tweet into 100 values, removes 
word order, and gives every recognized token equal influence.

The classifier comparison is much smaller than the representation comparison. 
Within TF-IDF, linear SVM exceeded Logistic Regression by only 0.0085 macro-F1. 
Within GloVe, Logistic Regression exceeded linear SVM by only 0.0017. This indicates 
that feature representation was the main source of performance difference in these experiments.

\subsection{Error Analysis}

The largest TF-IDF errors occur between adjacent intensity levels. 
Of 592 Extremely Negative tweets, 197 were predicted as Negative; of 599 Extremely 
Positive tweets, 181 were predicted as Positive. The model also predicted 214 of 1,041 Negative 
tweets as Positive. These errors are consistent with ambiguity in intensity, negation, quotation, 
sarcasm, and mixed opinions.

GloVe did not resolve these issues: its Logistic Regression model mapped 286 Extremely 
Negative tweets to Negative and 287 Extremely Positive tweets to Positive. Averaging pretrained 
word vectors removes compositional information that could distinguish ordinary from extreme 
sentiment. TF-IDF cannot fully model such context either, but it retains highly discriminative 
surface terms that mean pooling dilutes.

\subsection{Advantages and Limitations}

TF-IDF is sparse, task-specific, interpretable, and comparatively strong for this dataset, 
but it does not represent semantic similarity well and cannot preserve full word order. 
Pretrained GloVe provides dense semantic features and compact 100-dimensional arrays, 
but its static vectors assign the same representation to a word in every context, 
and averaging discards token order and intensity structure.

The dataset also imposes limitations. Labels separating ordinary from extreme sentiment 
may be subjective, and the tweets represent a specific pandemic period and topic. 
Results may not generalize to later language or other domains. Only a compact parameter grid 
was tested, so the results compare the implemented configurations rather than every possible 
classifier or embedding strategy.